% Complete actual-file proof audits: root and C, 20 November 2026.
\begin{lemma}[Proved: two-sided slopes for the saturated boundary grades]
\label{D:saturated-boundary-two-sided-slopes}
Use the fixed resolution in Proposition
\ref{D:tame-parabolic-boundedness}, refined once so that the
strict transforms $C_i$ of the original boundary are smooth
and pairwise disjoint, and their union with the exceptional
divisor has simple normal crossings without triple points.
Choose at every coefficient-Frobenius stage one minimal
extension $F_a$ of the canonical reflexive lattice $E_a$.
Extend the full parabolic flag from
$C_i-(C_i\cap\operatorname{Exc}(\pi))$ to $F_a|_{C_i}$
by saturation.  Denote its nonzero graded vector bundles by
$G_{a,i,j}$.

There is a constant $K$, depending only on the fixed
resolution, such that
\[
                -K\le\mu_{\min}(G_{a,i,j})
                    \le\mu_{\max}(G_{a,i,j})\le K
\]
for every rank, coefficient field, Kummer denominator,
stage, boundary component and grade.
\end{lemma}

\begin{proof}
Fix a positive exceptional divisor $Z$ with
$Z\cdot E_\nu<0$ for every exceptional component.
All minimal extensions have restrictions to $E_\nu$
in the same fixed two-sided slope windows.  The curve
Frobenius estimate, obtained from semistability of
Frobenius direct image and adjunction
\citep[Theorem~2.2]{C_Sun2008}
\citep[Proposition~1.3 and its proof]{C_MehtaPauly2007},
puts the restrictions of $F_Y^*F_a$ in fixed windows
as well.  Here $p$ is fixed, so the error in a single
Frobenius step is independent of rank.

On $Y-\operatorname{Exc}(\pi)$ the ordinary map
$\phi_a:F_Y^*F_a\to F_{a+1}$ is regular, and its inverse
has pole order at most $p-1$ along every $C_i$.
The bounded-pole lemma for exceptional slope windows
therefore gives one $B$, independent of all the
varying data, such that both maps extend as
\[
 \begin{split}
 F_Y^*F_a&\longrightarrow F_{a+1}(BZ),\\
 F_{a+1}&\longrightarrow
       F_Y^*F_a\bigl((p-1)\textstyle\sum_i C_i+BZ\bigr).
 \end{split}                                            \tag{D.SF.1}
\]
For the second application, twisting by the fixed
$(p-1)\sum_iC_i$ only shifts the exceptional windows
by fixed intersection numbers.  The bounded-pole
argument consequently applies to this target too.

The saturated flag on a smooth curve is a flag of
subbundles with locally free quotients.  At any point
of $C_i$ choose an adapted basis of its restriction;
that basis lifts to a basis of $F_a$ on a neighborhood
in the smooth surface.  Thus the full flag defines
a locally split parabolic bundle on the pair
$(Y,\sum_iC_i)$.  For a common denominator $n$ prime
to $p$, let $\mathcal F_a$ be the corresponding
bundle on the root stack
\[
                \mathcal X=\sqrt[n]{(Y,\sum_i C_i)}.
\]
Only the strict original boundary is rooted;
the exceptional divisor is not added to the logarithmic
boundary.

We claim that (D.SF.1) normalizes to maps
\[
 F_{\mathcal X}^*\mathcal F_a
        \longrightarrow\mathcal F_{a+1}(BZ),\qquad
 \mathcal F_{a+1}
        \longrightarrow F_{\mathcal X}^*\mathcal F_a(BZ).
                                                        \tag{D.SF.2}
\]
Away from $Z$ these are the genuine Frobenius
isomorphism of the tame local system and its inverse.
It remains to check that normalization does not
increase the exceptional pole order.

At a point of $C_i\cap Z$, take regular parameters
$t,z$, with $C_i=(t)$ and the exceptional component
$(z)$.  Write the coefficient of $BZ$ there as $b$.
On the Kummer chart $t=s^n$, the local ring is regular,
with parameters $s,z$.  A normalized root-matrix entry
is an ordinary matrix entry multiplied by an integral
power of $s$.  From (D.SF.1), after multiplication by
$z^b$, it can have poles only in $s$.  Its regularity
after inverting $z$ follows from the root-stack
Frobenius isomorphism away from the exceptional set.
Consequently it has nonnegative valuation at $(s)$.
Normality, or the intersection of the two localizations
at $s$ and $z$, now says that it is regular.
Exactly the same argument applies to the inverse:
the allowed ordinary pole $t^{-(p-1)}$ contributes
only an $s$-power and no $z$-pole.  This proves
(D.SF.2), with the same $B$, independently of $n$.
At points of $Z$ away from the strict boundary the
claim is simply (D.SF.1).

Let $N_i=\mathcal O_Y(C_i)|_{C_i}$ and let
$P_i=Z|_{C_i}$, an effective fixed divisor on $C_i$.
On the root gerbe over $C_i$, choose the character-line
convention in which its tautological line $\mathcal L$
satisfies $\mathcal L^n=N_i$ and the character-$j$
summand of $\mathcal F_a$ is
$G_{a,i,j}\otimes\mathcal L^j$.
If $pj=b+nc_j$, with $0\le b<n$, then
$0\le c_j\le p-1$.  Since multiplication by $p$ permutes
the characters, (D.SF.2) restricts, character by
character, to generically invertible maps
\[
 \begin{split}
 F_{C_i}^*G_{a,i,j}\otimes N_i^{c_j}
        &\longrightarrow G_{a+1,i,b}(BP_i),\\
 G_{a+1,i,b}
        &\longrightarrow
        F_{C_i}^*G_{a,i,j}\otimes N_i^{c_j}(BP_i).
 \end{split}                                            \tag{D.SF.3}
\]
Both are injective, because they are generically
isomorphisms between vector bundles on a smooth
integral curve.  This step uses the saturated flags
inside the same $F_a|_{C_i}$; it does not choose an
independent extension for each grade.

Put $K_i^0=B\deg P_i+(p-1)|\deg N_i|$.
The first injection gives
\[
 p\,\mu_{\max}(G_{a,i,j})
              \le\mu_{\max}(G_{a+1,i,b})+K_i^0.
\]
Dualizing the second injection gives precisely
the analogous inequality for the dual bundles:
\[
 p\,\mu_{\max}(G_{a,i,j}^{\vee})
              \le\mu_{\max}(G_{a+1,i,b}^{\vee})+K_i^0.
\]
The pairs consisting of the coefficient stage and
the character form finite cycles.  Taking the
maximum in each of the two inequalities around
such a cycle bounds both maxima by $K_i^0/(p-1)$.
Since $\mu_{\min}(G)=-\mu_{\max}(G^\vee)$, taking
the maximum over the finitely many $C_i$ proves
the assertion.
\end{proof}

\begin{remark}[The remaining quantitative presentation step]
The slope windows would give one fixed twist making
every grade, every flag subbundle and every quotient
globally generated with vanishing first cohomology
on each $C_i$.  Their section spaces then have
dimensions bounded by a constant times the rank.
The entire flag can be encoded by one adapted basis
of one such section space, with quadratically many
matrix coordinates.  A full bounded-degree
presentation of the evaluation, quotient and
Frobenius compatibility conditions, and the
Artinian openness argument for its unit classes,
must still be written before a quadratic-rank
representation bound is claimed.
\end{remark}
